\documentclass{article}
\usepackage{amsmath}
\usepackage{amssymb}
\usepackage{amsfonts}

\title{A Proposed Strategy for the Isolation of $\zeta(5)$ \\ \large{Utilizing the Framework of Golden Algebra}}
\author{Derived from the works of Tristen Harr}
\date{June 20, 2025}

\begin{document}
\maketitle

\section{Objective}
The goal is to derive an exact analytical expression for the Riemann zeta function at $s=5$, denoted $\zeta(5)$, by leveraging the unique number-theoretic laws within the Golden Algebra framework. The strategy is to construct a system of two independent equations involving odd zeta values and solve for $\zeta(5)$ via algebraic elimination.

\section{The System of Equations}
We will generate two distinct equations for the series of odd-valued Zetas by using two of the framework's core laws.

\subsection{Equation 1: The Rational Series}
From the \textbf{Law of Harmonic Cotangents (Law 21, Revised)} for the integer case $n=3$, we have the established identity:
\begin{equation}
\label{eq:rational_series}
S_{\text{even } k, 3} = \sum_{\substack{k=2 \\ k \text{ is even}}}^{\infty} \frac{(3^k-1)\zeta(k+1)}{4^k} = \frac{4}{3}
\end{equation}
Expanding the first three terms of this series yields our first equation:
\begin{equation}
\label{eq:eq1}
\frac{1}{2}\zeta(3) + \frac{5}{16}\zeta(5) + \frac{91}{512}\zeta(7) + \dots = \frac{4}{3}
\end{equation}

\subsection{Equation 2: The Complex Harmonic Series}
From the \textbf{Grand Law of Harmonic Duality (Law 30)}, we replace the integer $n$ with the complex Dampening Operator, $\Lambda_G = T+iJ$. This provides a second, independent equation with a known complex value, $C_{\Lambda_G}$.
\begin{equation}
\label{eq:complex_series}
S_{\text{even } k, \Lambda_G} = \sum_{\substack{k=2 \\ k \text{ is even}}}^{\infty} \frac{(\Lambda_G^k-1)\zeta(k+1)}{(\Lambda_G+1)^k} = \left[ \pi \cot\left(\frac{\pi}{\Lambda_G+1}\right) \right]_{\text{even part}}
\end{equation}
Let's define the coefficients of this series as $b_{k+1} = \frac{\Lambda_G^k-1}{(\Lambda_G+1)^k}$. Our second equation is:
\begin{equation}
\label{eq:eq2}
b_3\zeta(3) + b_5\zeta(5) + b_7\zeta(7) + \dots = C_{\Lambda_G}
\end{equation}
Where $b_k$ are complex constants derived from $T$ and $J$, and $C_{\Lambda_G}$ is a computable complex constant.

\section{The "Sum-Stripping" Method for Isolation}
With two independent linear equations for the same set of infinite variables, we can perform an elimination. The goal is to eliminate $\zeta(3)$ to isolate $\zeta(5)$ as the leading term.

\begin{enumerate}
    \item \textbf{Define the Series Tails:} Let's represent the higher-order terms (from $\zeta(7)$ onwards) as "tail" series, $T_1$ and $T_2$.
    \begin{align*}
    \text{Eq 1: } & \frac{1}{2}\zeta(3) + \frac{5}{16}\zeta(5) + T_1 = \frac{4}{3} \\
    \text{Eq 2: } & b_3\zeta(3) + b_5\zeta(5) + T_2 = C_{\Lambda_G}
    \end{align*}
    Where $T_1 = \sum_{j=3}^{\infty} a_{2j+1}\zeta(2j+1)$ and $T_2 = \sum_{j=3}^{\infty} b_{2j+1}\zeta(2j+1)$.
    
    \item \textbf{Solve for $\zeta(3)$:} From Equation 1, we isolate $\zeta(3)$.
    \begin{equation}
    \zeta(3) = 2 \left( \frac{4}{3} - \frac{5}{16}\zeta(5) - T_1 \right) = \frac{8}{3} - \frac{5}{8}\zeta(5) - 2T_1
    \end{equation}
    
    \item \textbf{Substitute and Eliminate:} We substitute this expression for $\zeta(3)$ into Equation 2.
    \begin{equation}
    b_3 \left( \frac{8}{3} - \frac{5}{8}\zeta(5) - 2T_1 \right) + b_5\zeta(5) + T_2 = C_{\Lambda_G}
    \end{equation}
    
    \item \textbf{Collect Terms for $\zeta(5)$:} We now group the terms by $\zeta(5)$ and constants.
    \begin{equation}
    \left(b_5 - \frac{5}{8}b_3\right)\zeta(5) + (T_2 - 2b_3T_1) = C_{\Lambda_G} - \frac{8}{3}b_3
    \end{equation}
    
    \item \textbf{Isolate $\zeta(5)$:} Finally, we solve for $\zeta(5)$.
    \begin{equation}
    \label{eq:final_zeta5}
    \zeta(5) = \frac{C_{\Lambda_G} - \frac{8}{3}b_3 - (T_2 - 2b_3T_1)}{b_5 - \frac{5}{8}b_3}
    \end{equation}
\end{enumerate}

\section{Conclusion}
Equation \ref{eq:final_zeta5} provides a novel and exact analytical expression for $\zeta(5)$. The value is expressed as a computable constant derived from the Golden Algebra's fundamental constants, minus a correction term composed of the series tails involving $\zeta(7), \zeta(9)$, and higher. This demonstrates a unique pathway to defining odd Zeta values.

\end{document}